\documentclass[12pt,a4paper]{article}
\usepackage[margin=2.5cm]{geometry}
\usepackage{amsmath,amssymb,amsfonts}
\usepackage{physics}
\usepackage{enumitem}
\usepackage{fancyhdr}
\usepackage{lastpage}
\usepackage{graphicx}
\usepackage{multicol}
\usepackage{array}
\usepackage{tabularx}
\usepackage{tikz}

% Page style
\pagestyle{fancy}
\fancyhf{}
\renewcommand{\headrulewidth}{0.4pt}
\renewcommand{\footrulewidth}{0.4pt}
\lhead{PHYS10012}
\rhead{Mock Examination}
\cfoot{Page \thepage\ of \pageref{LastPage}}

% Question numbering
\newcounter{questioncounter}
\newcommand{\question}[1]{%
  \stepcounter{questioncounter}%
  \vspace{0.5cm}%
  \noindent\textbf{Question \thequestioncounter.} \hfill [#1 marks]%
  \vspace{0.2cm}%
  \par%
}

% Sub-part environments
\newlist{parts}{enumerate}{2}
\setlist[parts,1]{label=(\alph*),leftmargin=2em,itemsep=0.3cm}
\setlist[parts,2]{label=(\roman*),leftmargin=2em,itemsep=0.2cm}

% Mark allocation in sub-parts
\renewcommand{\marks}[1]{\hfill [#1]}

% MCQ environment
\newcounter{mcqcounter}
\newcommand{\mcq}[1]{%
  \stepcounter{mcqcounter}%
  \vspace{0.3cm}%
  \noindent\textbf{\themcqcounter.} #1 \hfill [2 marks]%
  \vspace{0.1cm}%
}
\newlist{choices}{enumerate}{1}
\setlist[choices]{label=\Alph*.,leftmargin=2.5em,itemsep=0.1cm}

% Dirac notation shortcuts
\renewcommand{\ket}[1]{\left|#1\right\rangle}
\renewcommand{\bra}[1]{\left\langle #1\right|}
\renewcommand{\braket}[2]{\left\langle #1\middle|#2\right\rangle}
\renewcommand{\ketbra}[2]{\left|#1\right\rangle\!\left\langle #2\right|}
\renewcommand{\expect}[1]{\left\langle #1\right\rangle}

\begin{document}

% Title block
\begin{center}
{\Large\textbf{UNIVERSITY OF BRISTOL}}\\[0.3cm]
{\large Faculty of Science}\\[0.5cm]
{\Large\textbf{MOCK EXAMINATION --- Paper 6 (Easy)}}\\[0.3cm]
{\large\textbf{Core Physics I --- Classical, Quantum and Thermal Physics}}\\[0.2cm]
{\large Module Code: PHYS10012}\\[0.5cm]
{\large Duration: 3 Hours}\\[0.3cm]
{\large Total Marks: 100}\\[0.5cm]
\end{center}

\noindent\rule{\textwidth}{0.4pt}

\vspace{0.3cm}
\noindent\textbf{Instructions:}
\begin{itemize}[itemsep=0.1cm]
\item Answer \textbf{ALL} questions.
\item \textbf{Section A} (Oscillations and Waves): 10 multiple-choice questions, 20 marks total.
\item \textbf{Section B} (Quantum Physics): 10 multiple-choice questions, 20 marks total.
\item \textbf{Section C}: 4 long-answer questions, 60 marks total.
\item An approved calculator is permitted.
\item A formula sheet is provided at the end of this paper.
\item Show all working for Section C. Marks will be awarded for clear, logical solutions.
\end{itemize}

\noindent\rule{\textwidth}{0.4pt}
\newpage

% ============================================================
% SECTION A: OSCILLATIONS AND WAVES MCQ
% ============================================================
\section*{Section A: Oscillations and Waves}
\noindent\textit{Answer all 10 questions. Each question is worth 2 marks.}\\[0.3cm]

\setcounter{mcqcounter}{0}

% QUESTION A1: SHM definition
\mcq{A particle undergoes simple harmonic motion (SHM). Which of the following correctly describes the restoring force acting on the particle?}
\begin{choices}
\item The force is constant in magnitude and always directed toward the equilibrium position.
\item The force is proportional to the displacement from equilibrium and directed toward the equilibrium position: $F = -kx$.
\item The force is proportional to the square of the displacement from equilibrium.
\item The force is proportional to the velocity of the particle.
\end{choices}

% QUESTION A2: Period of mass-spring
\mcq{A mass of $m = 0.5\;\text{kg}$ is attached to a spring with spring constant $k = 200\;\text{N/m}$. What is the period of oscillation?}
\begin{choices}
\item $T = 0.157\;\text{s}$
\item $T = 0.201\;\text{s}$
\item $T = 0.314\;\text{s}$
\item $T = 0.628\;\text{s}$
\end{choices}

% QUESTION A3: SHM max velocity
\mcq{A mass oscillates in SHM with amplitude $A = 0.10\;\text{m}$ and angular frequency $\omega_0 = 5\;\text{rad/s}$. What is the maximum speed of the mass?}
\begin{choices}
\item $0.25\;\text{m/s}$
\item $0.50\;\text{m/s}$
\item $2.5\;\text{m/s}$
\item $5.0\;\text{m/s}$
\end{choices}

% QUESTION A4: Quality factor
\mcq{In a damped oscillating system, the quality factor $Q$ is defined as:}
\begin{choices}
\item $Q = \gamma/\omega_0$, the ratio of the damping constant to the natural frequency.
\item $Q = \omega_0/\gamma$, the ratio of the natural angular frequency to the damping constant.
\item $Q = \omega_0 \gamma$, the product of the natural frequency and damping constant.
\item $Q = \omega_d/\omega_0$, the ratio of the damped frequency to the natural frequency.
\end{choices}

% QUESTION A5: Wave speed
\mcq{A wave has frequency $f = 500\;\text{Hz}$ and wavelength $\lambda = 0.68\;\text{m}$. What is the speed of the wave?}
\begin{choices}
\item $170\;\text{m/s}$
\item $340\;\text{m/s}$
\item $500\;\text{m/s}$
\item $735\;\text{m/s}$
\end{choices}

% QUESTION A6: Standing waves
\mcq{A string of length $L$ is fixed at both ends and vibrates in its 3rd harmonic. How many antinodes are present on the string?}
\begin{choices}
\item $1$
\item $2$
\item $3$
\item $4$
\end{choices}

% QUESTION A7: Sound level
\mcq{The sound intensity level of $0\;\text{dB}$ corresponds to a sound intensity of:}
\begin{choices}
\item $0\;\text{W/m}^2$
\item $1\;\text{W/m}^2$
\item $10^{-12}\;\text{W/m}^2$
\item $10^{-6}\;\text{W/m}^2$
\end{choices}

% QUESTION A8: Doppler
\mcq{A sound source moves \emph{toward} a stationary observer at speed $v_s$ (where $v_s < v$, the speed of sound). The frequency heard by the observer is:}
\begin{choices}
\item Lower than the emitted frequency.
\item Equal to the emitted frequency.
\item Higher than the emitted frequency.
\item Zero, because the source is moving.
\end{choices}

% QUESTION A9: Beats
\mcq{Two tuning forks produce sound waves with frequencies $f_1 = 256\;\text{Hz}$ and $f_2 = 260\;\text{Hz}$. What is the beat frequency heard when both sound simultaneously?}
\begin{choices}
\item $2\;\text{Hz}$
\item $4\;\text{Hz}$
\item $258\;\text{Hz}$
\item $516\;\text{Hz}$
\end{choices}

% QUESTION A10: Simple pendulum
\mcq{The period of a simple pendulum undergoing small oscillations depends on which of the following quantities?}
\begin{choices}
\item The mass of the bob and the length of the string only.
\item The length of the string and the acceleration due to gravity only.
\item The amplitude of oscillation and the mass of the bob only.
\item The mass of the bob, the length of the string, and the acceleration due to gravity.
\end{choices}

\newpage

% ============================================================
% SECTION B: QUANTUM PHYSICS MCQ
% ============================================================
\section*{Section B: Quantum Physics}
\noindent\textit{Answer all 10 questions. Each question is worth 2 marks.}\\[0.3cm]

\setcounter{mcqcounter}{0}

% QUESTION B1: What is |psi>?
\mcq{In Dirac notation, the symbol $\ket{\psi}$ represents:}
\begin{choices}
\item A probability
\item A complex number
\item A quantum state (state vector/ket)
\item An operator
\end{choices}

% QUESTION B2: Inner product type
\mcq{In Dirac notation, the quantity $\braket{\phi}{\psi}$ is:}
\begin{choices}
\item A ket (column vector)
\item A bra (row vector)
\item An operator (matrix)
\item A complex scalar (complex number)
\end{choices}

% QUESTION B3: Born rule
\mcq{For two normalised quantum states $\ket{\phi}$ and $\ket{\psi}$, the quantity $|\braket{\phi}{\psi}|^2$ represents:}
\begin{choices}
\item The energy of the state $\ket{\psi}$
\item The probability of finding the system in state $\ket{\phi}$ given that it is in state $\ket{\psi}$
\item The expectation value of the Hamiltonian
\item The phase difference between the two states
\end{choices}

% QUESTION B4: Eigenvalue equation
\mcq{In the eigenvalue equation $\hat{A}\ket{a} = a\ket{a}$, the quantity $a$ is called:}
\begin{choices}
\item The expectation value
\item The normalisation constant
\item The eigenvalue
\item The eigenvector
\end{choices}

% QUESTION B5: Spin-1/2 S_z eigenvalues
\mcq{For a spin-$\frac{1}{2}$ particle, the possible results of measuring $S_z$ are:}
\begin{choices}
\item $0$ and $\hbar$
\item $+\hbar/2$ and $-\hbar/2$
\item $+\hbar$ and $-\hbar$
\item $+1$ and $-1$
\end{choices}

% QUESTION B6: Normalisation
\mcq{A spin-$\frac{1}{2}$ particle is in the state $\ket{\psi} = \alpha\ket{\!\uparrow} + \beta\ket{\!\downarrow}$, where $\alpha$ and $\beta$ are complex numbers. For this state to be normalised, which condition must hold?}
\begin{choices}
\item $\alpha + \beta = 1$
\item $\alpha^2 + \beta^2 = 1$
\item $|\alpha|^2 + |\beta|^2 = 1$
\item $|\alpha| + |\beta| = 1$
\end{choices}

% QUESTION B7: Hermitian
\mcq{An operator $\hat{A}$ is said to be Hermitian if:}
\begin{choices}
\item $\hat{A}^2 = \hat{I}$
\item $\hat{A}^\dagger = -\hat{A}$
\item $\hat{A}^\dagger = \hat{A}$
\item $\hat{A}^\dagger \hat{A} = \hat{I}$
\end{choices}

% QUESTION B8: Spin-1 states
\mcq{A particle with spin quantum number $s = 1$ has how many distinct spin states (i.e.\ distinct values of $m_s$)?}
\begin{choices}
\item $1$
\item $2$
\item $3$
\item $4$
\end{choices}

% QUESTION B9: Proton charge from quarks
\mcq{The proton has quark content $uud$. Given that the up quark has charge $+\frac{2}{3}e$ and the down quark has charge $-\frac{1}{3}e$, what is the total charge of the proton?}
\begin{choices}
\item $0$
\item $+\frac{1}{3}e$
\item $+e$
\item $+2e$
\end{choices}

% QUESTION B10: Electron classification
\mcq{An electron is classified as which type of particle?}
\begin{choices}
\item A boson and a baryon
\item A fermion and a lepton
\item A boson and a lepton
\item A fermion and a baryon
\end{choices}

\newpage

% ============================================================
% SECTION C: LONG ANSWER QUESTIONS
% ============================================================
\section*{Section C: Long Answer Questions}
\noindent\textit{Answer all 4 questions. Show all working.}

\setcounter{questioncounter}{0}

% ---- QUESTION C1: MECHANICS (15 marks) ----
\question{15}

\noindent\textbf{SUVAT, Forces, and Energy --- Direct Application}

\begin{parts}
\item A ball is thrown vertically upward from ground level with an initial speed of $20\;\text{m/s}$. Take $g = 9.81\;\text{m/s}^2$.
\begin{parts}
\item Find the time taken to reach the maximum height. \qmarks{1}
\item Find the maximum height reached. \qmarks{1}
\item Find the total time of flight (i.e.\ the time for the ball to return to ground level). \qmarks{1}
\end{parts}

\item A $10\;\text{kg}$ box is pushed across a horizontal floor by a horizontal force of $60\;\text{N}$. The coefficient of kinetic friction between the box and the floor is $\mu_k = 0.4$.
\begin{parts}
\item List all the forces acting on the box and describe the direction of each (i.e.\ describe the free body diagram). \qmarks{1}
\item Calculate the kinetic friction force. \qmarks{1}
\item Find the acceleration of the box. \qmarks{1}
\end{parts}

\item A $2\;\text{kg}$ ball is dropped from a height of $5\;\text{m}$ above the ground. It reaches the ground and then slides onto a frictionless horizontal surface where it encounters a spring with spring constant $k = 500\;\text{N/m}$. Assume no energy is lost.
\begin{parts}
\item What is the speed of the ball just before it reaches the ground? \qmarks{2}
\item Using energy conservation, find the maximum compression of the spring. \qmarks{2}
\end{parts}

\item A $3\;\text{kg}$ object moving at $4\;\text{m/s}$ to the right collides with a $1\;\text{kg}$ object that is initially at rest.
\begin{parts}
\item If the collision is perfectly inelastic (the objects stick together), find the final velocity of the combined mass. \qmarks{2}
\item Calculate the kinetic energy lost in the perfectly inelastic collision. \qmarks{1}
\item If instead the collision is elastic, use the elastic collision formula from the formula sheet to find the final velocity of the $1\;\text{kg}$ object. \qmarks{2}
\end{parts}
\end{parts}

\newpage

% ---- QUESTION C2: PROPERTIES OF MATTER (15 marks) ----
\question{15}

\noindent\textbf{Heat, Ideal Gas, and Basic Thermodynamics}

\begin{parts}
\item $0.5\;\text{kg}$ of water at $20\;^\circ\text{C}$ is heated to $100\;^\circ\text{C}$ and then completely converted to steam at $100\;^\circ\text{C}$. Calculate the total energy required. Use $c_\text{water} = 4200\;\text{J/(kg\,K)}$ and $L_\text{vap} = 2.26\times 10^6\;\text{J/kg}$. \qmarks{3}

\item A container holds $1\;\text{mol}$ of ideal monatomic gas at $300\;\text{K}$ and $1.0\times 10^5\;\text{Pa}$.
\begin{parts}
\item Find the volume of the gas. \qmarks{1}
\item The gas is heated at constant pressure until its temperature reaches $600\;\text{K}$. Find the new volume. \qmarks{1}
\item Calculate the work done by the gas during this isobaric expansion. \qmarks{1}
\item Calculate the heat absorbed by the gas. \qmarks{1}
\end{parts}

\item The gas from part~(b), now at $600\;\text{K}$ and pressure $1.0\times 10^5\;\text{Pa}$, undergoes an adiabatic expansion to twice its volume. Take $\gamma = 5/3$.
\begin{parts}
\item Find the final temperature after the adiabatic expansion. \qmarks{2}
\item Find the final pressure. Is it higher or lower than the initial pressure $1.0\times 10^5\;\text{Pa}$? \qmarks{2}
\end{parts}

\item A Carnot engine operates between a hot reservoir at $T_h = 500\;\text{K}$ and a cold reservoir at $T_c = 300\;\text{K}$. The engine absorbs $Q_h = 1000\;\text{J}$ of heat per cycle from the hot reservoir.
\begin{parts}
\item What is the efficiency of the engine? \qmarks{1}
\item How much work does the engine do per cycle? \qmarks{1}
\item How much heat is rejected to the cold reservoir per cycle? \qmarks{1}
\item What is the entropy change of the hot reservoir per cycle? Of the cold reservoir? Of the universe? \qmarks{1}
\end{parts}
\end{parts}

\newpage

% ---- QUESTION C3: OSCILLATIONS AND WAVES (15 marks) ----
\question{15}

\noindent\textbf{Direct SHM and Wave Applications}

\begin{parts}
\item A mass $m = 0.25\;\text{kg}$ is attached to a spring with spring constant $k = 100\;\text{N/m}$ and oscillates with amplitude $A = 4\;\text{cm} = 0.04\;\text{m}$.
\begin{parts}
\item Find the angular frequency $\omega_0$ and the period $T$. \qmarks{2}
\item Find the total energy of the oscillator. \qmarks{1}
\item Find the maximum speed of the mass. \qmarks{1}
\end{parts}

\item A transverse wave on a string is described by $y(x,t) = 0.04\sin(3x - 12t)\;\text{m}$, where $x$ is in metres and $t$ is in seconds.
\begin{parts}
\item Determine the amplitude, wavelength, frequency, and wave speed. \qmarks{2}
\item In which direction is the wave travelling? \qmarks{1}
\end{parts}

\item A string of length $L = 2\;\text{m}$ is fixed at both ends. The wave speed on the string is $v = 200\;\text{m/s}$.
\begin{parts}
\item Find the fundamental frequency $f_1$. \qmarks{1}
\item Find the frequencies of the 2nd and 3rd harmonics. \qmarks{1}
\item Sketch the displacement patterns for the first three harmonics on the string, clearly indicating nodes (N) and antinodes (A). \qmarks{2}
\end{parts}

\item A fire engine siren emits sound at $f = 800\;\text{Hz}$. The speed of sound is $v = 340\;\text{m/s}$.
\begin{parts}
\item The fire engine approaches a stationary observer at $v_s = 30\;\text{m/s}$. What frequency does the observer hear? \qmarks{2}
\item The fire engine passes and moves away from the observer at $v_s = 30\;\text{m/s}$. What frequency does the observer now hear? \qmarks{2}
\end{parts}
\end{parts}

\newpage

% ---- QUESTION C4: QUANTUM PHYSICS (15 marks) ----
\question{15}

\noindent\textbf{Basic Dirac Notation and Spin Measurements}

\begin{parts}
\item A spin-$\frac{1}{2}$ particle is in the state
\[
\ket{\psi} = \frac{3}{5}\ket{\!\uparrow_z} + \frac{4}{5}\ket{\!\downarrow_z}.
\]
\begin{parts}
\item Verify that the state is normalised. \qmarks{1}
\item If $S_z$ is measured, what are the possible outcomes and their probabilities? \qmarks{2}
\end{parts}

\item Using the spin eigenstates from the formula sheet, express $\ket{\!\uparrow_z}$ and $\ket{\!\downarrow_z}$ in terms of $\ket{\!\uparrow_x}$ and $\ket{\!\downarrow_x}$. Hence write $\ket{\psi}$ from part~(a) in the $x$-basis and find the probability of measuring $S_x = +\hbar/2$. \qmarks{4}

\item Consider the operator $\hat{A} = \ketbra{\!\uparrow}{\uparrow\!} - \ketbra{\!\downarrow}{\downarrow\!}$ (written in the $z$-basis).
\begin{parts}
\item Write $\hat{A}$ as a $2\times 2$ matrix in the $z$-basis. \qmarks{1}
\item What are the eigenvalues and eigenvectors of $\hat{A}$? \qmarks{2}
\item Is $\hat{A}$ Hermitian? How do you know? \qmarks{1}
\end{parts}

\item The state $\ket{\psi} = \frac{1}{\sqrt{2}}\bigl(\ket{\!\uparrow_z} + \ket{\!\downarrow_z}\bigr)$ is an eigenstate of which spin operator ($S_x$, $S_y$, or $S_z$)? What is the corresponding eigenvalue? If $S_z$ is measured on this state, what is the expectation value $\expect{S_z}$? \qmarks{4}
\end{parts}

\newpage

% ============================================================
% FORMULA SHEET
% ============================================================
% EQUATION SHEET — PHYS10012 Core Physics I
% This file is \input{} at the end of each exam paper

\newpage
\section*{Formula Sheet}
\noindent\rule{\textwidth}{0.4pt}

\subsection*{Mathematical Formulae}

\begin{align*}
&\text{Binomial expansion:} & (1+x)^n &\approx 1 + nx + \frac{n(n-1)}{2!}x^2 + \cdots \quad (|x| \ll 1) \\[4pt]
&\text{Euler's formula:} & e^{i\theta} &= \cos\theta + i\sin\theta \\[4pt]
&\text{Trig identities:} & \sin A + \sin B &= 2\cos\!\left(\tfrac{A-B}{2}\right)\sin\!\left(\tfrac{A+B}{2}\right) \\
& & \cos A + \cos B &= 2\cos\!\left(\tfrac{A-B}{2}\right)\cos\!\left(\tfrac{A+B}{2}\right)
\end{align*}

\subsection*{Mechanics}

\begin{align*}
&\text{SUVAT:} & v &= u + at, \quad s = ut + \tfrac{1}{2}at^2, \quad v^2 = u^2 + 2as, \quad s = \tfrac{1}{2}(u+v)t \\[4pt]
&\text{Dot product:} & \mathbf{a}\cdot\mathbf{b} &= |\mathbf{a}||\mathbf{b}|\cos\theta = a_x b_x + a_y b_y + a_z b_z \\[4pt]
&\text{Cross product:} & \mathbf{a}\times\mathbf{b} &= |\mathbf{a}||\mathbf{b}|\sin\theta\;\hat{\mathbf{n}} = \begin{vmatrix}\hat{\mathbf{i}}&\hat{\mathbf{j}}&\hat{\mathbf{k}}\\a_x&a_y&a_z\\b_x&b_y&b_z\end{vmatrix} \\[4pt]
&\text{Rocket equation:} & \Delta v &= u \ln\!\left(\frac{m_0}{m_f}\right) \\[4pt]
&\text{Gravitational PE:} & U &= -\frac{GMm}{r} \\[4pt]
&\text{Escape velocity:} & v_e &= \sqrt{\frac{2GM}{R}} \\[4pt]
&\text{Force from PE:} & F &= -\frac{dU}{dx} \\[4pt]
&\text{Elastic collision (1D):} & v_1' &= \frac{m_1 - m_2}{m_1 + m_2}v_1 + \frac{2m_2}{m_1+m_2}v_2 \\[4pt]
&\text{Centre of mass:} & \mathbf{R}_\text{cm} &= \frac{\sum m_i\mathbf{r}_i}{\sum m_i} \\[4pt]
&\text{Reduced mass:} & \mu &= \frac{m_1 m_2}{m_1 + m_2} \\[4pt]
&\text{Centripetal accel:} & a_c &= \frac{v^2}{r} = \omega^2 r \\[4pt]
&\text{Parallel axis thm:} & I &= I_\text{cm} + Md^2 \\[4pt]
&\text{Atwood machine:} & a &= \frac{(m_1-m_2)g}{m_1+m_2}, \quad T = \frac{2m_1 m_2 g}{m_1+m_2}
\end{align*}

\noindent\textbf{Moments of inertia} (about axis through centre of mass):
\begin{center}
\begin{tabular}{ll}
Thin ring (radius $R$): & $I = MR^2$ \\
Solid disk (radius $R$): & $I = \tfrac{1}{2}MR^2$ \\
Solid sphere (radius $R$): & $I = \tfrac{2}{5}MR^2$ \\
Hollow sphere (radius $R$): & $I = \tfrac{2}{3}MR^2$ \\
Thin rod (length $L$, about centre): & $I = \tfrac{1}{12}ML^2$ \\
Thin rod (length $L$, about end): & $I = \tfrac{1}{3}ML^2$ \\
\end{tabular}
\end{center}

\subsection*{Properties of Matter}

\begin{align*}
&\text{Ideal gas:} & PV &= nRT \\[4pt]
&\text{Heat capacities:} & C_P - C_V &= R \quad\text{(per mole)} \\[4pt]
&\text{Adiabatic:} & PV^\gamma &= \text{const}, \quad TV^{\gamma-1} = \text{const} \\[4pt]
&\text{Isothermal work:} & W &= nRT\ln\!\left(\frac{V_2}{V_1}\right) \\[4pt]
&\text{Carnot efficiency:} & \eta &= 1 - \frac{T_c}{T_h} \\[4pt]
&\text{Entropy:} & dS &= \frac{dQ_\text{rev}}{T}, \quad \Delta S = mc\ln\!\left(\frac{T_2}{T_1}\right), \quad \Delta S = nR\ln\!\left(\frac{V_2}{V_1}\right) \\[4pt]
&\text{Lennard-Jones:} & V(r) &= 4\varepsilon\left[\left(\frac{a}{r}\right)^{12} - \left(\frac{a}{r}\right)^{6}\right] \\[4pt]
&\text{Mean KE:} & \langle KE \rangle &= \tfrac{3}{2}k_B T \\[4pt]
&\text{MB speeds:} & c_\text{mp} &= \sqrt{\frac{2k_BT}{m}}, \quad c_\text{avg} = \sqrt{\frac{8k_BT}{\pi m}}, \quad c_\text{rms} = \sqrt{\frac{3k_BT}{m}}
\end{align*}

\subsection*{Oscillations and Waves}

\begin{align*}
&\text{SHM:} & x(t) &= A\cos(\omega_0 t + \phi), \quad \omega_0 = \sqrt{k/m} \\[4pt]
&\text{Damped SHM:} & x(t) &= Ae^{-\gamma t/2}\cos(\omega_d t + \phi), \quad \omega_d = \sqrt{\omega_0^2 - \gamma^2/4} \\[4pt]
&\text{Quality factor:} & Q &= \omega_0/\gamma \\[4pt]
&\text{Forced amplitude:} & A(\omega) &= \frac{F_0/m}{\sqrt{(\omega_0^2 - \omega^2)^2 + \gamma^2\omega^2}} \\[4pt]
&\text{Phase lag:} & \tan\delta &= \frac{\gamma\omega}{\omega_0^2 - \omega^2} \\[4pt]
&\text{Pendulums:} & T_\text{simple} &= 2\pi\sqrt{L/g}, \quad T_\text{physical} = 2\pi\sqrt{I/(Mgh)} \\[4pt]
&\text{Wave equation:} & \frac{\partial^2 y}{\partial x^2} &= \frac{1}{v^2}\frac{\partial^2 y}{\partial t^2} \\[4pt]
&\text{String speed:} & v &= \sqrt{T/\mu} \\[4pt]
&\text{Sound in gas:} & v &= \sqrt{\gamma P/\rho} \\[4pt]
&\text{Harmonic wave:} & y &= A\sin(kx - \omega t) \\[4pt]
&\text{String impedance:} & Z &= \mu v = \sqrt{\mu T} \\[4pt]
&\text{Reflection:} & r &= \frac{Z_1 - Z_2}{Z_1 + Z_2}, \quad t = \frac{2Z_1}{Z_1+Z_2} \\[4pt]
&\text{Power (string):} & \langle P\rangle &= \tfrac{1}{2}\mu\omega^2 A^2 v \\[4pt]
&\text{Decibels:} & \beta &= 10\log_{10}(I/I_0), \quad I_0 = 10^{-12}\;\text{W/m}^2 \\[4pt]
&\text{Doppler:} & f' &= f\frac{v \pm v_o}{v \mp v_s} \\[4pt]
&\text{Group velocity:} & v_g &= \frac{d\omega}{dk} \\[4pt]
&\text{String harmonics:} & f_n &= \frac{nv}{2L} \quad (n=1,2,3,\ldots) \\[4pt]
&\text{Open-closed pipe:} & f_n &= \frac{nv}{4L} \quad (n=1,3,5,\ldots)
\end{align*}

\subsection*{Quantum Physics}

\begin{align*}
&\text{Schr\"{o}dinger eq:} & i\hbar\frac{d}{dt}\ket{\psi(t)} &= H\ket{\psi(t)} \\[4pt]
&\text{Time evolution:} & \ket{\psi(t)} &= \sum_n c_n e^{-iE_n t/\hbar}\ket{E_n}, \quad U(t) = e^{-iHt/\hbar} \\[4pt]
&\text{Spin operators:} & S_z &= \frac{\hbar}{2}\bigl(\ketbra{\!\uparrow}{\uparrow\!} - \ketbra{\!\downarrow}{\downarrow\!}\bigr) \\
& & S_x &= \frac{\hbar}{2}\bigl(\ketbra{\!\uparrow}{\downarrow\!} + \ketbra{\!\downarrow}{\uparrow\!}\bigr) \\
& & S_y &= \frac{\hbar}{2}\bigl(-i\ketbra{\!\uparrow}{\downarrow\!} + i\ketbra{\!\downarrow}{\uparrow\!}\bigr) \\[4pt]
&\text{Spin eigenstates:} & \ket{\!\uparrow_x} &= \tfrac{1}{\sqrt{2}}\bigl(\ket{\!\uparrow_z} + \ket{\!\downarrow_z}\bigr), \quad \ket{\!\downarrow_x} = \tfrac{1}{\sqrt{2}}\bigl(\ket{\!\uparrow_z} - \ket{\!\downarrow_z}\bigr) \\
& & \ket{\!\uparrow_y} &= \tfrac{1}{\sqrt{2}}\bigl(\ket{\!\uparrow_z} + i\ket{\!\downarrow_z}\bigr), \quad \ket{\!\downarrow_y} = \tfrac{1}{\sqrt{2}}\bigl(\ket{\!\uparrow_z} - i\ket{\!\downarrow_z}\bigr) \\[4pt]
&\text{Pauli matrices:} & \sigma_x &= \begin{pmatrix}0&1\\1&0\end{pmatrix}, \quad \sigma_y = \begin{pmatrix}0&-i\\i&0\end{pmatrix}, \quad \sigma_z = \begin{pmatrix}1&0\\0&-1\end{pmatrix}
\end{align*}

\subsection*{Physical Constants}

\begin{center}
\begin{tabular}{lll}
Speed of light & $c$ & $3.00 \times 10^{8}$ m/s \\
Planck's constant & $h$ & $6.626 \times 10^{-34}$ J$\cdot$s \\
Reduced Planck's constant & $\hbar$ & $1.055 \times 10^{-34}$ J$\cdot$s \\
Boltzmann constant & $k_B$ & $1.381 \times 10^{-23}$ J/K \\
Gas constant & $R$ & $8.314$ J/(mol$\cdot$K) \\
Avogadro's number & $N_A$ & $6.022 \times 10^{23}$ mol$^{-1}$ \\
Gravitational constant & $G$ & $6.674 \times 10^{-11}$ N$\cdot$m$^2$/kg$^2$ \\
Acceleration due to gravity & $g$ & $9.81$ m/s$^2$ \\
Electron mass & $m_e$ & $9.109 \times 10^{-31}$ kg \\
Proton mass & $m_p$ & $1.673 \times 10^{-27}$ kg \\
Elementary charge & $e$ & $1.602 \times 10^{-19}$ C \\
Mass of Earth & $M_\oplus$ & $5.97 \times 10^{24}$ kg \\
Radius of Earth & $R_\oplus$ & $6.37 \times 10^{6}$ m \\
\end{tabular}
\end{center}

\noindent\rule{\textwidth}{0.4pt}


\end{document}
